\documentclass[../../../include/open-logic-section]{subfiles}

\begin{document}

\olfileid{sth}{replacement}{intro}
\olsection{ভূমিকা}

প্রতিস্থাপন সেই স্বতঃসিদ্ধ-ছক, যা $\ZF$ ও
~$\Z$-এর মধ্যে পার্থক্য গড়ে দেয়।
\crefrange{sth:ordinals::chap}{sth:spine::chap}
জুড়ে আমরা এটি অবলম্বন করেছি। এই অধ্যায়ে
অবশেষে প্রশ্নটি বিবেচনা করব: প্রতিস্থাপনের
ন্যায্যতা আছে কি?

প্রশ্নটির গুরুত্ব বোঝার জন্য লক্ষ করা দরকার,
প্রতিস্থাপন আসলে বেশ \emph{শক্তিশালী}।
এই অধ্যায়ে (এবং আবার
\olref[sth][card-arithmetic][fix]{sec}-এ)
এর শক্তি কতখানি, তার ধারণা পাব। তাতেই
বোঝা যাবে, এর ন্যায্যতা প্রতিষ্ঠা করা
সত্যিই প্রয়োজন।

দুই ধরনের ন্যায্যতা নিয়ে আলোচনা করব। মোটামুটি,
কোনো স্বতঃসিদ্ধের ফলের ভিত্তিতে তাকে ন্যায্য
প্রতিপন্ন করার চেষ্টা হলো \emph{বাহ্যিক} ন্যায্যতা;
সংশ্লিষ্ট গাণিতিক ধারণাগুলিই স্বতঃসিদ্ধটিকে
সমর্থন করে—এমন প্রস্তাবের মাধ্যমে তাকে ন্যায্য
প্রতিপন্ন করার চেষ্টা হলো \emph{অন্তর্নিহিত}
ন্যায্যতা। এই অধ্যায়ে পার্থক্যটি আরও স্পষ্ট
হবে, যদিও এটি বৃহত্তর আলোচনার কেবল সূচনা।
আরও জানতে বিশেষ করে
\citeauthor{Maddy1988a}-এর লেখা দেখো
(\citeyear{Maddy1988a} এবং
\citeyear{Maddy1988b})।

\end{document}
